\documentclass[11pt]{article}
\usepackage[a4paper,margin=1in]{geometry}
\usepackage{amsmath,amssymb,amsthm}
\usepackage{booktabs}
\usepackage{graphicx}
\usepackage{hyperref}
\usepackage{enumitem}
\usepackage{algorithm}
\usepackage{algpseudocode}
\title{CS-768 (Learning with Graphs) Course Project Report\\\large GraphTransLWG: Hybrid GNN--Transformer for Long-Range Graph Reasoning}
\author{Eshaan Pareek (22B3986) \and Tejas Sharma (22B0909) \and Aditya Vema Reddy Kesari (22B3985) \and Matam Kushaal (23B1290)}
\date{}

\begin{document}
\maketitle

\section{Introduction}
Graph Neural Networks (GNNs) aggregate neighborhood information using message passing, but depth increases can cause oversmoothing and oversquashing. This limits long-range dependency modeling. Our project studies a GraphTrans-style hybrid architecture that interleaves local graph convolutions with global self-attention to capture both local topology and distant interactions.

\section{Problem Setup}
Given a graph $G=(V,E)$ with node features $X\in\mathbb{R}^{|V|\times d_x}$ and graph label $y$, learn $f_\theta(G)\to\hat{y}$ for graph-level prediction.

\paragraph{Message Passing (local bias).}
A generic layer updates node $i$ as:
\[
\mathbf{h}^{(\ell+1)}_i = \phi\left(\mathbf{h}^{(\ell)}_i,\; \square_{j\in\mathcal{N}(i)} \psi(\mathbf{h}^{(\ell)}_i,\mathbf{h}^{(\ell)}_j,\mathbf{e}_{ij})\right),
\]
where $\square$ is an aggregation operator.

\paragraph{Global attention (long-range context).}
With hidden matrix $H\in\mathbb{R}^{n\times d}$:
\[
Q=HW_Q,\quad K=HW_K,\quad V=HW_V,
\]
\[
\mathrm{Attn}(H)=\mathrm{softmax}\left(\frac{QK^\top}{\sqrt{d_h}}\right)V.
\]
This enables every node to interact with every other node in one layer.

\section{Architecture}
The implemented model follows:
\begin{enumerate}[leftmargin=1.5em]
    \item Input projection: $X\mapsto H^{(0)}\in\mathbb{R}^{n\times d}$.
    \item Add a graph-level \texttt{[CLS]} token per graph for pooling.
    \item Repeat for $L$ blocks:
    \begin{itemize}
        \item GNN sublayer for local structure encoding.
        \item Multi-head self-attention sublayer for global context.
        \item Position-wise MLP sublayer.
        \item Residual paths and normalization in transformer stack.
    \end{itemize}
    \item Readout from \texttt{[CLS]} embedding and linear head for logits.
\end{enumerate}

\section{Implementation Mapping to Repository}
\begin{itemize}
    \item \texttt{models/full\_model.py}: \texttt{GraphTransConfig}, \texttt{GraphTransModel}, \texttt{ModelTrainConfig}, CLS pooling path.
    \item \texttt{models/transformer.py}: transformer block composition.
    \item \texttt{models/gnn.py}: local message-passing layers.
    \item \texttt{models/attention.py}: attention mechanism implementation.
    \item \texttt{models/train.py}: training loop with optimizer/loss/backprop.
    \item \texttt{data\_utils/config\_objects.py}: dataset and training configuration objects.
    \item \texttt{data\_utils/dataset\_utils.py}: dataset loading and preprocessing utilities.
    \item \texttt{trainers/base\_trainer.py}: generic trainer with grad clipping + scheduler hooks.
    \item \texttt{trainers/flag\_trainer.py}: FLAG adversarial perturbation training.
\end{itemize}

\section{Training Objective and Optimization}
For classification, with logits $z_g=f_\theta(G_g)$ and labels $y_g$, we minimize:
\[
\mathcal{L}(\theta)=\frac{1}{B}\sum_{g=1}^{B} \mathrm{CE}(z_g,y_g).
\]
Training uses Adam-family optimization with optional gradient clipping:
\[
\theta \leftarrow \theta - \eta\,\nabla_\theta \mathcal{L}.
\]

\subsection{FLAG Robust Training}
FLAG injects adversarial feature perturbations $\delta$ in hidden/node input space:
\[
\delta^{t+1} = \Pi_{\epsilon}\left(\delta^t + \alpha\,\mathrm{sign}(\nabla_\delta \mathcal{L}(\theta,\delta^t))\right).
\]
Model parameters are updated against perturbed examples, improving robustness and regularization.

\section{Complexity Discussion}
Let $n=|V|$, hidden width $d$, edges $m=|E|$.
\begin{itemize}
    \item GNN layer cost is typically $\mathcal{O}(m d)$.
    \item Full attention cost is $\mathcal{O}(n^2 d)$ in time and $\mathcal{O}(n^2)$ in memory for attention scores.
\end{itemize}
Hence, hybrid architectures trade better long-range modeling for higher cost on large graphs.

\section{Benchmarks and Experimental Scope}
Configuration files indicate evaluation on OGBG-Code2, OGBG-MolPCBA, NCI1, and NCI109 with multiple variants (GCN/GIN/PNA with and without virtual node settings).

\section{Conclusion}
This project establishes a technical and reproducible foundation for graph representation learning with GraphTrans-style hybrids. The key benefit is combining local inductive bias from message passing with global receptive fields from attention.

\end{document}